\documentclass{article}
\usepackage[UTF8]{ctex}
\usepackage{amsmath,amssymb,amsfonts}
\usepackage{enumitem}
\usepackage{geometry}
\geometry{a4paper, margin=1in}
% === 错题本排版优化（不参与题目标识） ===
\setlength{\parindent}{0pt}
\setlength{\parskip}{0.35em plus 0.10em minus 0.08em}
\setlength{\emergencystretch}{3em}
\setlist{itemsep=0.25em, topsep=0.40em, parsep=0pt}
\allowdisplaybreaks[2]
% === 错题本排版优化结束 ===
\begin{document}
\section*{陕西师范大学2026年数学分析真题}
\section*{189．陕西师范大学2026年数学分析试卷}

\subsection{题目1}
1．求函数 \(\frac{1}{x^{2}+3 x+2}\) 在 \(x=-4\) 处的幂级数展开式和收敛域。

\subsection{题目2}
2．设函数 \(y=y(x)\) 由参数方程 \(\left\{\begin{array}{l}x=2 e^{t}+t+1 \\ y=4(t-1) e^{t}+t^{2}\end{array}\right.\) 确定，计算 \(\frac{\mathrm{d}^{2} y}{\mathrm{~d} x^{2}}\) ．

\subsection{题目3}
3．设 \(f(x)\) 的一个原函数为 \(F(x)>0\) ，且 \(F(0)=1\) ，当 \(x \geq 0\) 时，有 \(f(x) F(x)=\sin ^{2} 2 x\) ，求 \(f(x)\) ．

\subsection{题目4}
4．计算曲线积分

\[
\int_{L}\left[2\left(x^{2}-x\right) e^{2 x}-x y\right] \mathrm{d} x+\left[x^{2}-(y+2) e^{y}\right] \mathrm{d} y
\]

式中 \(L\) 是从 \((0,0)\) 经曲线 \(y=x^{2}-2 x\) 到点 \(A(4,8)\) 的一段弧。

\subsection{题目5}
5．计算曲面积分

\[
\iint_{\sum} x \mathrm{~d} y \mathrm{~d} z+y \mathrm{~d} z \mathrm{~d} x+z \mathrm{~d} x \mathrm{~d} y .
\]

其中 \(\sum\) 为上半球面 \(z=\sqrt{R^{2}-x^{2}-y^{2}}(R>0)\) 的上侧．

\subsection{题目6}
6．解答如下问题：\\
（1）讨论函数 \(f(x)=(x-1)(x-2) D(x)\) 的连续点和间断点，并判断间断点的类型，其中 \(D(x)\) 是狄利克雷函数

\[
D(x)= \begin{cases}1, & x \text { 为有理数; } \\ 0, & x \text { 为无理数. }\end{cases}
\]

（2）给出实数集 \(\mathbb{R}\) 上只有 3 个连续点的函数．

\subsection{题目7}
7．利用交换积分次序计算 \(\int_{1}^{2} \mathrm{~d} x \int_{\sqrt{x}}^{x} \sin \frac{\pi x}{2 y} \mathrm{~d} y+\int_{2}^{4} \mathrm{~d} x \int_{\sqrt{x}}^{2} \sin \frac{\pi x}{2 y} \mathrm{~d} y\) ．

\subsection{题目8}
8．讨论反常积分 \(I=\int_{0}^{\frac{\pi}{2}} \ln \cos x \mathrm{~d} x\) 的敛散性，若收敛，计算 \(I\) 的值．

\subsection{题目9}
9．设 \(f(x)\) 在 \([0,1]\) 上连续，在 \((0,1)\) 内可导，且 \(\int_{0}^{1} f(x) \mathrm{d} x=0\) ，记 \(F(x)=\int_{0}^{x} x f(t) \mathrm{d} t\) ．

\begin{itemize}
  \item[（1）] 求 \(F^{\prime}(x)\) 。
  \item[（2）] 证明：在 \((0,1)\) 内至少存在一点 \(\xi\) ，使得 \(\int_{0}^{\xi} f(x) \mathrm{d} x=-\xi f(\xi)\) ．
  \item[（3）] 证明：在 \((0,1)\) 内至少存在一点 \(x_{0}\) ，使得 \(2 f\left(x_{0}\right)+x_{0} f^{\prime}\left(x_{0}\right)=0\) ．
\end{itemize}

\subsection{题目10}
10．设 \(a_{n}>0, b_{n}>0\) ，证明：

\begin{itemize}
  \item[（1）] 若存在 \(\alpha>0\) ，使得 \(\frac{b_{n}}{b_{n+1}} a_{n}-a_{n+1} \geq \alpha(n=1,2, \cdots)\) ，则级数 \(\sum_{n=1}^{\infty} b_{n}\) 收敛。
  \item[（2）] 若 \(\sum_{n=1}^{\infty} \frac{1}{a_{n}}\) 发散，且 \(\frac{b_{n}}{b_{n+1}} a_{n}-a_{n+1} \leq 0(n=1,2, \cdots)\) ，则级数 \(\sum_{n=1}^{\infty} b_{n}\) 发散．
\end{itemize}

\subsection{题目11}
11．设 \(x \geq 0, n\) 为正整数，\(f(x)=\int_{0}^{x}\left(t-t^{2}\right) \sin ^{2 n} t \mathrm{~d} t\) ，证明：\(f(x) \leq \frac{1}{(2 n+2)(2 n+3)}\) ．

\subsection{题目12}
12．设 \(f(x)\) 是 \([0,1]\) 上的正值连续函数，且 \(f(x)\) 单调递减，证明：

\[
\frac{\int_{0}^{1} x f^{2}(x) \mathrm{d} x}{\int_{0}^{1} x f(x) \mathrm{d} x} \leq \frac{\int_{0}^{1} f^{2}(x) \mathrm{d} x}{\int_{0}^{1} f(x) \mathrm{d} x}
\]

\subsection{题目13}
13．设 \(S(x)=\sum_{n=1}^{\infty} n e^{-n x}\) ，其中 \(x>0\) ，求证：

\begin{itemize}
  \item[（1）] 函数项级数 \(\sum_{n=1}^{\infty} n e^{-n x}\) 在 \((0,+\infty)\) 上不一致收敛。
  \item[（2）] 和函数 \(S(x)\) 在 \((0,+\infty)\) 上连续可导。
\end{itemize}

\subsection{题目14}
14．证明：若在区间 \(I\) 上函数列 \(\left\{f_{n}(x)\right\}\) 一致收敛，且对每个正整数 \(n, f_{n}(x)\) 都在区间 \(I\) 上有界，则函数列 \(\left\{f_{n}(x)\right\}\) 在 \(I\) 上一致有界．
\end{document}
